% T08 source: Краткие термины и обозначения.tex
This section provides a concise conceptual overview of the formulation and its
key notation; formal definitions and deterministic assumptions are introduced
in the \emph{Formal Model} section. The model describes a fixed software-development workload that is completed to
a predefined acceptance criterion. The no-AI baseline consists of sequential
completion of the same workload by a single developer. A complete glossary of
terms and the full notation system are provided in
Appendix~\ref{app:glossary}.

\begin{description}[style=nextline]
  \item[Task and mode.] Task $i$ has size $Z_i$ and is performed in operating
  mode $r$: interactive, delegated, or fully autonomous. The mode determines
  the sequence of task specification, agent execution, review, and correction
  of the result. The main M4 analysis concerns the interactive and delegated
  modes; the fully autonomous mode, which has no mandatory human-attention
  phase, serves as a limiting case.

  \item[Normalized AI-assisted duration.] The quantity
  $k_i^{(r)}=T_i^{ai}(r)/(Z_iX)$ relates the duration required to obtain an
  accepted result to the baseline time for the same task. It is specific to
  the ``task--tool--mode'' triplet, not to the AI tool in general. This is a
  descriptive normalization; a causal interpretation requires a separate
  identification design.

  \item[Phases.] The normalized duration is decomposed as
  $k_i^{(r)}=a_i^{(r)}+h_i^{(r)}$. Autonomous agent phases require an agent
  slot; human-attention phases simultaneously require the developer's
  attention and retention of the assigned slot.

  \item[Configured parallelism.] The parameter $P$ denotes the number of agent
  streams configured in the given configuration and available to the schedule;
  coordination overhead is evaluated at this value. Additional unused capacity
  under an unchanged configuration does not itself generate overhead.

  \item[Configuration.] The notation $\sigma=(P,r)$ or $\sigma=(P,\rho)$ is
  shorthand when the AI tool and acceptance rules are fixed. If the tool or the
  acceptance criterion for an accepted result changes, the change must be
  stated explicitly in the comparison, although these elements are not added
  as dimensions of $\sigma$.

  \item[Schedule and makespan.] A feasible schedule respects the task-DAG
  dependencies, phase order, the unit capacity of the human resource, and local
  slot retention. Its makespan is denoted by $T(\pi)$, and the minimum makespan
  over feasible schedules by $T^*$.
\end{description}

\begin{center}
\small
\begin{tabularx}{\textwidth}{@{}>{$}l<{$}>{\raggedright\arraybackslash}X@{}}
\toprule
\text{Notation} & \text{Meaning} \\
\midrule
X & baseline effort per unit of work without AI \\
Z_i & size of task $i$ in units of $X$ \\
k_i=a_i+h_i & normalized duration and its agent and human components \\
P & configured number of agent streams in the configuration \\
G=(V,E) & task-dependency DAG \\
W=X\sum_iZ_ik_i & total work when $C=\gamma=1$ \\
A,\ H & total autonomous agent work and human work \\
C(P),\ \gamma(P) & multipliers for cross-stream integration overhead and human attention-switching overhead \\
T(\pi),\ T^* & makespan of a particular schedule and the optimal makespan \\
\bottomrule
\end{tabularx}
\end{center}
